\subsection{Evaluación incidencial de la coordenatización terminal}
\label{subsec:k-incidencias-residentes}

El estado pleno conserva orientación, hoja, acarreo, firma de frontera,
memoria y libro dodecafásico. Sobre su subred compatible, el lector firmado
es la proyección de la coordenada generada y coincide con
\[
 R_{\mathrm{sgn}}=\operatorname{pr}_{U}
 =\Pi_H\circ\operatorname{pr}_{C}
 =\Pi_H\circ\mathcal E_{108}^{90,120}\circ\mathcal R_{12}.
\]
Estas identidades operan sobre el estado enriquecido, no sobre la proyección
cruda de los ciento ocho pasos que olvida aquellos campos. Un descenso por
esa proyección requeriría constancia en sus fibras y no se presupone aquí.

Para explicitar la presentación incidencial, sea \(\mathscr L\) el libro
orientado del estado pleno. Cada visita conserva ruta, posición APP, hoja,
instante, multiplicidad e incidencia de frontera. La evaluación en la hoja
de suma o producto determina \(n_v=r_v+9q_v\), con \(1\leq r_v\leq9\);
el residuo no es una entrada independiente. Para
\(E_m=\{9(m-1)+1,\ldots,9m\}\), se definen
\[
 \mathscr V_m=\{v:t(v)\in E_m\},\qquad
 \mathscr T_m=\{\theta:\operatorname{origen}(\theta)\in E_m\},
\]
y los recuentos
\[
\begin{aligned}
 A_m&=\sum_{v\in\mathscr V_m}w(v)\mathbf1_{\{1,3,5,7\}}(r(v)),\\
 C_m&=-\sum_{\theta\in\mathscr T_m}u(\theta)\sigma(\theta)
     =\sum_{j\geq0}\bigl(\nu^-_{120,m}(j)-\nu^+_{120,m}(j)\bigr),\\
 V_m&=\sum_{v\in\mathscr V_m}w(v)\mathbf1_{\{9\}}(r(v))
                         \epsilon_\partial(v).
\end{aligned}
\]
Aquí \(w\) y \(u\) son las multiplicidades de visitas y trazas,
\(\sigma\) conserva su orientación y \(\epsilon_\partial\) vale \(1\) en
la corona y \(-1\) en el interior. Las trazas de \(\nu^\pm\) se enumeran
mediante la regla del libro; la suma entera requiere soporte finito o
una construcción que determine su significado. El número de incidencias
cronológicas y la longitud reducida \(120(1+3j)\) son magnitudes distintas:
su relación conserva la unidad de longitud del lector.

Con \([x]_{10}\in\{0,\ldots,9\}\), se mantienen también los cocientes
de \(x=[x]_{10}+10\lfloor x/10\rfloor\). La lectura decimal y la
extracción transversal son
\[
\begin{aligned}
 z_m&=100[A_m]_{10}+10[C_m]_{10}+[V_m]_{10},\\
 \mathcal E_{\mathrm{cnt}}(\mathscr L)
   &=\bigl((z_m-z_{m+3})_{m=1}^{12},
           (z_m-z_{m+4})_{m=1}^{12},\sum_{m=1}^{12}z_m\bigr).
\end{aligned}
\]
Los índices son cíclicos módulo doce. Así, \(\mathcal E_{\mathrm{cnt}}\)
es la presentación \((D_3z,D_4z,Q(z))\) del extractor, con la misma
convención que en \S\ref{subsec:k-transversal}. Sus argumentos son los
recuentos del libro; no son \(K\), \(U\), alfa ni coordenadas transversales
esperadas. La presentación por recuentos se aplica cuando se despliega
el libro completo. No afirma que éste sea recuperable después de olvidar
orientación y memoria. La evaluación armónica y la inversión orbital
desarrolladas a continuación conservan el dominio compatible de la
coordenatización terminal.
